\subsection{Block dynamics on cells and a coarse-scale hypothesis in four
dimensions}\label{rC12:coarse}

A route to the averaged condition $\mathrm{BP}(n,\varpi)$ of
Definition~\ref{def:avgcube} is to coarse-grain the lattice into cells of
side $d$, to treat the inside of the cells by the Efron--Stein inequality,
and to treat the coarse variables by comparison with $\beta=0$. This
subsection carries out that route in four dimensions and locates the step
where it stops. The results are the following.
\begin{enumerate}[label=\textup{(\alph*)},leftmargin=*]
\item A finite-size criterion: a gap of the block dynamics on patches of
$m^4$ cells, uniform in the boundary condition, gives
$\mathrm{BP}(2d-1,\varpi)$ on every torus with $\varpi$ depending on that gap
and on $m$ only (Theorem~\ref{rC12:finite}). Blocks that share no plaquette
give commuting projections, and this replaces the disjointness used in the
criteria of Knabe~\cite{Kna88} and of Gosset and Mozgunov~\cite{GM16}.
\item A coarse-scale hypothesis, a density of the coarse links close to the
Haar density uniformly in the boundary condition
(Definition~\ref{rC12:C}), holds for every compact $G$, every weight and
every coupling once $d$ is large, with an error exponentially small in $d$
(Theorem~\ref{rC12:segment}). Links on one straight line share no
plaquette, so a coarse link is a product of $d-2$ conditionally independent
variables. The hypothesis therefore holds in the Coulomb phase of the
$\U(1)$ Villain theory, where the averaged condition fails for box measures,
and it implies neither the averaged condition nor confinement
(Proposition~\ref{rC12:coulomb}).
\item For every nonconstant weight two adjacent cells are not
conditionally independent given the coarse links, nor given any functions
of the two cells that leave out one pair of links of a plaquette crossing
their common face (Proposition~\ref{rC12:nogo}). With the quantitative
substitute $(\mathrm T_{d,m}(\kappa))$ of Definition~\ref{rC12:T}, an
approximate tensorisation of the variance over the cells given the coarse
links, the route closes and gives $\mathrm{BP}(2d-1,\varpi)$ with
$\varpi=1/(1024\kappa+16384)$ (Theorem~\ref{rC12:reduction}).
\item At fixed $d$ and small coupling the patch gap follows from comparison
with $\beta=0$; this gives the averaged condition in an explicit range
without Dobrushin coefficients or gauge fixing
(Corollary~\ref{rC12:strong}), a range far smaller than that of
Corollary~\ref{cor:bpranges}.
\end{enumerate}
The hypothesis $(\mathrm T_{d,m}(\kappa))$ and the patch gap are not proved
at weak coupling. Nothing in this subsection proves
$\mathrm{BP}(n_k,\varpi)$ along the continuum tuning, condition $D(n,
\alpha)$, condition (E), an area law or the mass gap.

\emph{Setting.} As in Sec.~\ref{ssec:averaged}, $G\subset\U(N)$ is compact
and nontrivial, $\varphi$ is a positive continuous class function on $G$,
$\Lambda=\Lambda_L=(\Z/L\Z)^4$, $E$ and $P$ are the sets of links and
plaquettes, $\mu$ is the measure with weight $\varphi$, and $h$ denotes
Haar probability measure (on $G$, or the product of Haar measures on $G^A$
for a set $A$ of links). Put
\[
\omega_\varphi=\log\frac{\max\varphi}{\min\varphi}.
\]
For $\SU(2)$ with $\varphi=e^{2\beta\ReTr(\cdot)}$, $\omega_\varphi=8\lvert
\beta\rvert=2\lvert\beta_W\rvert$. For a set $A$ of links, $P_A$ is the
conditional expectation under $\mu$ given the links outside $A$; the
conditional law of $U_A$ has density proportional to $\prod_{p\cap A\ne
\emptyset}\varphi(U_p)$ with respect to $h$ on $G^A$ and depends only on the
links that share a plaquette with a link of $A$. Two sets $A,A'$ of links
are \emph{non-interacting} if no plaquette contains a link of $A$ and a link
of $A'$. A plaquette with base point $y$ in the plane $\{\mu,\nu\}$ consists
of links with base points $y$, $y+e_\mu$, $y+e_\nu$, so the base points of
its links differ by at most one in each coordinate; a link lies on six
plaquettes, and two distinct links share at most one plaquette.

Fix an integer $d\ge2$ dividing $L$, put $K=L/d$, $n=2d-1$, and fix a shift
$a\in[0,d)^4$. For $j$ in the coarse torus $T_K=(\Z/K\Z)^4$ let $r_j=a+dj$
be the \emph{root} of $j$, let the \emph{block} $B_j=V_{r_j}$ be the set of
links with base point in $r_j+[0,n)^4$, and write $\mathcal P_j=P_{B_j}$,
which is the cube move $\mathcal P_{r_j}$ of Sec.~\ref{ssec:averaged} for
cubes of side $n$. Displacements $v=j'-j$ in $T_K$ are represented in
$(-K/2,K/2]^4$.

\subsubsection{A finite-size criterion}

\begin{lemma}\label{rC12:commute}
\begin{enumerate}[label=\textup{(\roman*)},leftmargin=*]
\item For non-interacting sets $A,A'\subset E$ of links, $P_AP_{A'}=
P_{A'}P_A=P_{A\cup A'}$, and $(I-P_A)(I-P_{A'})$ is an orthogonal projection.
\item If $K\ge4$ and $\lvert j'-j\rvert_\infty\ge2$, then $B_j$ and $B_{j'}$
are non-interacting.
\item The blocks $B_j$, $j\in T_K$, cover $E$.
\item If $f\in L^2(\mu)$ and $\mathcal P_jf=f$ for every $j\in T_K$, then $f$
is almost surely constant.
\end{enumerate}
\end{lemma}

\begin{proof}
(i) Every link lies on a plaquette, so $A\cap A'=\emptyset$, and no link of
$A'$ shares a plaquette with a link of $A$. Given the links outside $A\cup
A'$, the law of $(U_A,U_{A'})$ has density proportional to $\prod_{p\cap(A
\cup A')\ne\emptyset}\varphi(U_p)$, and each factor involves the links of
exactly one of $A$, $A'$; so the density is a product of a function of
$U_A$ and a function of $U_{A'}$, with coefficients depending on the links
outside $A\cup A'$. Hence $U_A$ and $U_{A'}$ are conditionally independent,
the conditional law $\gamma_A$ of $U_A$ given the links outside $A$ does not
depend on $U_{A'}$, and likewise for $A'$. Therefore $P_AP_{A'}f$ and
$P_{A'}P_Af$ are both the integral of $f$ against $\gamma_A\otimes
\gamma_{A'}$, which is $P_{A\cup A'}f$. Then $(I-P_A)(I-P_{A'})=I-P_A-P_{A'}+
P_{A\cup A'}$ is self-adjoint and idempotent.

(ii) Let $\lvert v_i\rvert\ge2$ for $v=j'-j$. In the coordinate $i$ the base
points of the links of $B_j$ and $B_{j'}$ lie in $dj_i+a_i+[0,2d-2]$ and
$dj_i+a_i+dv_i+[0,2d-2]$ modulo $L$. If $v_i\ge2$, their distance in $\Z/L\Z$
is at least $\min\{(v_i-2)d+2,(K-v_i-2)d+2\}\ge2$, since $v_i\le K/2$ and
$K\ge4$; $v_i\le-2$ is symmetric. A plaquette meeting both would have two
links whose base points differ by at least two in the coordinate $i$.

(iii) The sets $r_j+[0,d)^4$ partition $\Lambda_L$, and $[0,d)\subset[0,n)$.

(iv) The density $Z^{-1}\prod_p\varphi(U_p)$ of $\mu$ with respect to $h$ on
$G^E$ is bounded above and below by positive constants, so $\mu$ and $h$ have
the same null sets and $L^2(\mu)=L^2(h)$. If $\mathcal P_jf=f$, then $f$ is
almost surely equal to a function of $U_{E\setminus B_j}$, and so it is fixed
by the conditional expectation $Q_j$ of $h$ given $U_{E\setminus B_j}$, which
integrates out $U_{B_j}$. By Fubini's theorem the $Q_j$ commute and their
product integrates out the links of $\bigcup_jB_j=E$, so $f=\int f\,dh$
almost surely.
\end{proof}

The next lemma is a finite-size criterion of the kind of
Knabe~\cite{Kna88} and of Gosset and Mozgunov~\cite{GM16}. In four
dimensions the displacements $v$ with $\lvert v\rvert_\infty=1$ are $80$
in number, and those with $\lvert v\rvert_\infty\ge2$ give commuting
projections by Lemma~\ref{rC12:commute}. For $s\in T_K$ and $m\ge4$ let
$S_s=s+[0,m-1]^4$ be the \emph{patch} at $s$ and $S^\circ_s=s+[0,m-2]^4$
the set of its $(m-1)^4$ block indices; put
\[
H=\sum_{j\in T_K}(I-\mathcal P_j),\qquad H_s=\sum_{j\in S^\circ_s}(I-
\mathcal P_j).
\]

\begin{lemma}\label{rC12:knabe}
Let $m\ge4$, $p=m-1$ and $K\ge2m$, and suppose that $H_s^2\ge\gamma H_s$
for every $s\in T_K$, with $\gamma>0$. Then $H^2\ge\varepsilon H$ with
\[
\varepsilon=1-\frac{(1-\gamma)p^4+72p^3-168p^2+96p-16}{p^3(p-2)}\ge
\frac{\gamma p-74}{p-2}.
\]
In particular $\varepsilon\ge\gamma/2$ if $m-1\ge148/\gamma$.
\end{lemma}

\begin{proof}
Write $\mathbf P_j=I-\mathcal P_j$ and $A_v=\sum_j\mathbf P_j\mathbf P_{j+v}$.
Since $K\ge2m$, $j$ and $j+v$ lie in a common $S^\circ_s$ only if $\lvert v_i
\rvert\le m-2$ for all $i$, and then for exactly $c_v=\prod_{i=1}^4(p-\lvert
v_i\rvert)$ points $s$. Let $O$ be the set of the $80$ displacements with
$\lvert v\rvert_\infty=1$, $D$ the set of those with $\lvert v\rvert_\infty
\ge2$ and $D'$ the set of $v\in D$ with all $\lvert v_i\rvert\le m-2$. For
$v\in D$, $\mathbf P_j\mathbf P_{j+v}$ is an orthogonal projection by
Lemma~\ref{rC12:commute}(i),(ii), so $A_v\ge0$. Summing over $s$,
\[
\sum_sH_s^2=c_0H+\sum_{v\in O}c_vA_v+\sum_{v\in D'}c_vA_v\ge\gamma\sum_sH_s
=\gamma c_0H .
\]
Put $c_*=p^3(p-2)$; then $c_v\le c_*$ for $v\in D'$, since one factor of
$c_v$ is at most $p-2$ and the others at most $p$. For projections
$\mathbf P,\mathbf P'$, $(\mathbf P\pm\mathbf P')^2\ge0$ gives $-(\mathbf P+
\mathbf P')\le\mathbf P\mathbf P'+\mathbf P'\mathbf P\le\mathbf P+\mathbf
P'$, so $-2H\le A_v+A_{-v}\le2H$; $O$ is symmetric and $c_{-v}=c_v$, hence
\[
\sum_{v\in O}(c_*-c_v)A_v\ge-\Sigma H,\qquad\Sigma=\sum_{v\in O}\lvert c_*-
c_v\rvert,
\]
and
\[
\begin{split}
c_*\sum_{v\in O}A_v&=\sum_{v\in O}c_vA_v+\sum_{v\in O}(c_*-c_v)A_v\\
&\ge-\bigl((1-\gamma)c_0+\Sigma\bigr)H-\sum_{v\in D'}c_vA_v .
\end{split}
\]
Since $H^2=\sum_vA_v=H+\sum_{O}A_v+\sum_DA_v$,
\[
\begin{split}
H^2&\ge\Bigl(1-\frac{(1-\gamma)c_0+\Sigma}{c_*}\Bigr)H\\
&\quad+\sum_{v\in D'}\Bigl(1-\frac{c_v}{c_*}\Bigr)A_v+\sum_{v\in D
\setminus D'}A_v,
\end{split}
\]
and the last two sums are nonnegative. It remains to compute $\Sigma$. A
$v\in O$ with $k$ nonzero coordinates has $c_v=p^{4-k}(p-1)^k$, and there are
$8$, $24$, $32$, $16$ of them for $k=1,2,3,4$. The differences $c_v-c_*$ are
$p^3$, $p^2$, $-p(p^2-3p+1)$ and $-2p(p-1)(p-2)+1$, with the signs shown
since $p\ge3$, so
\[
\begin{split}
\Sigma&=8p^3+24p^2+32p(p^2-3p+1)\\
&\quad+16\bigl(2p(p-1)(p-2)-1\bigr),
\end{split}
\]
which is $72p^3-168p^2+96p-16$. With $c_0=p^4$ this is the formula for
$\varepsilon$. As $-168p^2+96p-16<0$, $\varepsilon\ge1-((1-\gamma)p+72)/(p-
2)=(\gamma p-74)/(p-2)$, and if $\gamma p\ge148$ this is at least $(\gamma
p-74)/p\ge\gamma/2$.
\end{proof}

For a patch $S_s$ let $\Pi_s=\bigcup_{j\in S^\circ_s}B_j$, the set of links
with base point in $r_s+[0,md-1)^4$. For $\omega\in G^{E\setminus\Pi_s}$ let
$\nu^\omega_s$ be the probability measure on $G^{\Pi_s}$ with density
proportional to $\prod_{p\cap\Pi_s\ne\emptyset}\varphi(U_p)$, $U$ equal to
$\omega$ off $\Pi_s$; it is the conditional law of $U_{\Pi_s}$ under $\mu$
given $U_{E\setminus\Pi_s}=\omega$. For $j\in S^\circ_s$ let
$\mathcal P^\omega_j$ be the conditional expectation under $\nu^\omega_s$
given the links of $\Pi_s\setminus B_j$.

\begin{definition}\label{rC12:patch}
The \emph{patch gap} $\gamma_{d,m}$ is the largest $\gamma\ge0$ such that
\[
\gamma\Var_{\nu^\omega_s}(g)\le\sum_{j\in S^\circ_s}\nu^\omega_s\bigl((g-
\mathcal P^\omega_jg)^2\bigr)
\]
for every $s$, every $\omega$ and every real $g\in L^2(\nu^\omega_s)$.
\end{definition}

The law $\nu^\omega_s$ involves only the plaquettes meeting $\Pi_s$, whose
links have base points in a cube of side $md+1$; for $L\ge md+2$ these are
the links and plaquettes of $\Z^4$, and translations carry the patches of
one shift to those of another. So $\gamma_{d,m}$ depends on $G$, $\varphi$,
$d$ and $m$ only. It is the gap of the dynamics that resamples the blocks of
one patch, uniform in the configuration outside the patch.

\begin{theorem}\label{rC12:finite}
Let $d\ge2$, $m\ge4$, $d\mid L$ and $L\ge2md$, and suppose that
$\gamma_{d,m}\ge\gamma>0$ and that $\varepsilon$, defined from $\gamma$ and
$m$ as in Lemma~\ref{rC12:knabe}, is positive. Then every real $f\in
L^2(\mu)$ satisfies
\[
\frac{\varepsilon}{16}\Var_\mu(f)\le\frac1{n^4}\sum_{x\in\Lambda_L}\mu\bigl(
(f-\mathcal P_xf)^2\bigr),\qquad n=2d-1.
\]
For even $L$ this is $\mathrm{BP}(2d-1,\varepsilon/16)$, and
$\mathrm{BP}_{\rm loc}(2d-1,\varepsilon/16)$ and, for $G=\SU(N)$,
$\mathrm{BP}_0(2d-1,\varepsilon/16)$ hold as well. If $m-1\ge148/\gamma$,
then $\varepsilon/16\ge\gamma/32$.
\end{theorem}

\begin{proof}
Fix $a$. Since $L\ge2md$, $K\ge2m\ge8$. First, $H_s^2\ge\gamma H_s$ for
every $s$. Write $U=(\sigma,\omega)$ with $\sigma\in G^{\Pi_s}$, $\omega\in
G^{E\setminus\Pi_s}$; under $\mu$ the conditional law of $\sigma$ given
$\omega$ is $\nu^\omega_s$. For $j\in S^\circ_s$, $B_j\subset\Pi_s$, and the
conditional law of $U_{B_j}$ given the other links is the same under $\mu$
and under $\nu^\omega_s$; so $(\mathcal P_jF)(\cdot,\omega)=\mathcal
P^\omega_j(F(\cdot,\omega))$ for $F\in L^2(\mu)$ and almost every $\omega$,
and $H_s$ acts on each fibre as $H^\omega_s=\sum_{j\in S^\circ_s}(I-\mathcal
P^\omega_j)$. This operator is self-adjoint and nonnegative, it annihilates
the constants and so preserves their orthogonal complement, and on that
complement $\ip g{H^\omega_sg}=\sum_j\nu^\omega_s((g-\mathcal P^\omega_j
g)^2)\ge\gamma\lVert g\rVert^2$ by Definition~\ref{rC12:patch}. Hence
the spectrum of $H^\omega_s$ lies in $\{0\}\cup[\gamma,\infty)$ and
$(H^\omega_s)^2\ge\gamma H^\omega_s$; integrating over $\omega$,
$\lVert H_sF\rVert^2\ge\gamma\ip F{H_sF}$.

By Lemma~\ref{rC12:knabe}, $H^2\ge\varepsilon H$, so the spectrum of $H$
lies in $\{0\}\cup[\varepsilon,\infty)$. If $HF=0$ then $\sum_j\lVert F-
\mathcal P_jF\rVert^2=\ip F{HF}=0$, so $F$ is constant by
Lemma~\ref{rC12:commute}(iv). Hence for $g=f-\mu(f)$, which is orthogonal
to the kernel,
\[
\varepsilon\Var_\mu(f)\le\ip g{Hg}=\sum_{j\in T_K}\mu\bigl((f-\mathcal
P_{r_j}f)^2\bigr).
\]
As $a$ runs over $[0,d)^4$ and $j$ over $T_K$, $r_j=a+dj$ runs over
$\Lambda_L$ once. Averaging over $a$ gives $\varepsilon d^4\Var_\mu(f)\le
\sum_x\mu((f-\mathcal P_xf)^2)$, and $d^4\ge n^4/16$. For even $L$,
$L\ge8d\ge2n+4$, so the inequality is $\mathrm{BP}(2d-1,\varepsilon/16)$
(on all real $f$, in particular on the gauge-invariant ones); the local
and neutral forms follow by Theorem~\ref{thm:bplocal}(i), and the last claim
by Lemma~\ref{rC12:knabe}.
\end{proof}

The theorem reduces the averaged condition on every large torus to a
finite-volume statement, uniform in the boundary condition, about patches of
side $md$. The passage costs a factor: at $\beta=0$, $\gamma_{d,m}\ge1$ by
Lemma~\ref{rC12:anova} below, and the theorem gives
$\mathrm{BP}(2d-1,\varepsilon/16)$ with $\varepsilon<1$, against
$\mathrm{BP}(n,1)$ from Proposition~\ref{prop:avgsup}(ii).

\subsubsection{The reference measure and the coarse links}

The following decomposition is that of Hoeffding~\cite{Hoe48} and of the
proof of the Efron--Stein inequality~\cite{ES81}, and it is the one used in
Proposition~\ref{prop:avgsup}(ii).

\begin{lemma}\label{rC12:anova}
Let $A$ be a finite set, $(X_e)_{e\in A}$ independent random variables with
values in measurable spaces, $\rho$ their joint law, and for $B\subset A$
let $Q_B$ be the conditional expectation given $X_{A\setminus B}$.
\begin{enumerate}[label=\textup{(\roman*)},leftmargin=*]
\item If $B_1,\dots,B_k$ cover $A$, then $\sum_{i}\rho((g-Q_{B_i}g)^2)\ge
\Var_\rho(g)$ for every real $g\in L^2(\rho)$.
\item If $A$ is partitioned into $C_1,\dots,C_k$, then $\Var_\rho(g)\le
\sum_i\rho((g-Q_{C_i}g)^2)$ for every real $g\in L^2(\rho)$.
\end{enumerate}
In particular, for $\varphi\equiv1$, $\gamma_{d,m}\ge1$ for all $d$ and
$m$.
\end{lemma}

\begin{proof}
Write $Q_e=Q_{\{e\}}$. By Fubini's theorem $Q_B$ integrates out $X_B$, the
$Q_e$ are commuting orthogonal projections, and $Q_B=\prod_{e\in B}Q_e$.
Expanding $I=\prod_{e\in A}(Q_e+(I-Q_e))$ gives $I=\sum_{F\subset A}\Delta_F$
with $\Delta_F=\prod_{e\notin F}Q_e\prod_{e\in F}(I-Q_e)$. Each $\Delta_F$ is
an orthogonal projection, $\Delta_F\Delta_{F'}=0$ for $F\ne F'$ (a factor
$Q_e(I-Q_e)$ occurs for $e$ in the symmetric difference), $\Delta_\emptyset
g=\rho(g)$, and $Q_B\Delta_F$ equals $\Delta_F$ if $F\cap B=\emptyset$ and $0$
otherwise. Hence, with $g_F=\Delta_Fg$,
\[
\rho\bigl((g-Q_Bg)^2\bigr)=\sum_{F\cap B\ne\emptyset}\lVert g_F\rVert^2,
\qquad\Var_\rho(g)=\sum_{F\ne\emptyset}\lVert g_F\rVert^2 .
\]
In (i) every $F\ne\emptyset$ meets some $B_i$, and (ii) is the special case
of a partition. For $\varphi\equiv1$, $\nu^\omega_s$ is the product of Haar
measures on $G^{\Pi_s}$, the blocks $B_j$, $j\in S^\circ_s$, cover $\Pi_s$,
and $\mathcal P^\omega_j=Q_{B_j}$; (i) gives $\gamma_{d,m}\ge1$.
\end{proof}

Part (ii) is the Efron--Stein inequality. Now the coarse variables. For $j\in
T_K$ and a direction $\nu$, the \emph{segment} $c=(j,\nu)$ is the set of the
$d$ links $(r_j+te_\nu,\nu)$, $0\le t<d$, which runs from $r_j$ to
$r_{j+e_\nu}$; its \emph{coarse link} is the ordered product
\[
W_c=U_{r_j,\nu}U_{r_j+e_\nu,\nu}\cdots U_{r_j+(d-1)e_\nu,\nu},
\]
and its \emph{middle links} are those with $1\le t\le d-2$, a set
$\mathrm{mid}(c)$. A gauge transformation $g$ acts by $W_{(j,\nu)}\mapsto
g_{r_j}W_{(j,\nu)}g_{r_{j+e_\nu}}^{-1}$, so the $W_c$ form a lattice gauge
field on $T_K$; its plaquettes are the holonomies of the boundaries of the
$d\times d$ squares spanned by the roots. The \emph{cell} of $j$ is the set
$D_j$ of links with base point in $r_j+[0,d)^4$; the cells partition $E$,
and the four segments $(j,\nu)$ lie in $D_j$. Lines of links are the sets
$\{(y+te_\nu,\nu):t\in\Z/L\Z\}$, and the segments lie on the \emph{coarse
lines}, those through roots; a link of direction $\nu$ lies on a coarse line
exactly when the coordinates of its base point other than $\nu$ are
congruent to those of $a$ modulo $d$.

\begin{lemma}\label{rC12:lines}
\begin{enumerate}[label=\textup{(\roman*)},leftmargin=*]
\item No plaquette contains two links of the same line.
\item Let $d\ge3$. A plaquette that contains a middle link of a segment
contains no other link of any coarse line.
\item Let $M\subset E$ be a set of links no two of which share a plaquette.
Given the links outside $M$, the variables $U_e$, $e\in M$, are independent,
and $U_e$ has a density $h_e$ with respect to $h$ satisfying
$e^{-6\omega_\varphi}\le h_e\le e^{6\omega_\varphi}$.
\item For $\varphi\equiv1$ the coarse links are independent and Haar
distributed, and the $U_{D_j}$, $j\in T_K$, are independent given all the
coarse links.
\end{enumerate}
\end{lemma}

\begin{proof}
(i) The two links of a plaquette in the plane $\{\mu,\nu\}$ that have
direction $\nu$ have base points $y$ and $y+e_\mu$, which differ in a
coordinate other than $\nu$, so they lie on different lines.

(ii) Let $e=(r_j+te_\nu,\nu)$ with $1\le t\le d-2$, and let the plaquette lie
in the plane $\{\nu,\mu\}$. Its other link of direction $\nu$ has base point
$r_j+te_\nu\pm e_\mu$, whose coordinate $\mu$ is $a_\mu\pm1$ modulo $d$, not
$a_\mu$ since $d\ge3$, so it is on no coarse line. Its links of direction
$\mu$ have base points $r_j+te_\nu$, $r_j+(t+1)e_\nu$ or these minus $e_\mu$,
whose coordinate $\nu$ is $a_\nu+t$ or $a_\nu+t+1$ with $1\le t<t+1\le d-1$,
so they are on no coarse line either.

(iii) The conditional density of $U_M$ is proportional to $\prod_{p\cap M\ne
\emptyset}\varphi(U_p)$, and each such $p$ contains exactly one link of $M$;
so it is a product over $e\in M$ of functions of $U_e$, the one for $e$ being
the product over the six plaquettes containing $e$. Normalised to
mean one, $h_e$ has maximum at least one and minimum at most one, and the
ratio of the two is at most $e^{6\omega_\varphi}$.

(iv) The segments are disjoint sets of independent Haar variables, and the
product of $d$ independent Haar variables is Haar distributed. The $U_{D_j}$
are independent and each $W_{(j,\nu)}$ is a function of $U_{D_j}$, so the
pairs $(U_{D_j},(W_{(j,\nu)})_\nu)$ are independent, and conditioning on all
second components keeps the first ones independent.
\end{proof}

By (iv), at $\beta=0$ the coarse links are coordinates in which the reference
law is a product: the Bianchi identity of the coarse plaquettes is carried
by the parametrisation, not by the measure. By
Lemma~\ref{rC12:anova}(ii), the $\beta=0$ dynamics that resamples one
coarse link at a time has gap one on all functions of the coarse links, in
particular on the coarse gauge-invariant functions.

\subsubsection{The coarse-scale hypothesis holds at every coupling}

For a patch $S_s$ let $C_s$ be the set of the segments $(j,\nu)$ with $j$ and
$j+e_\nu$ in $S_s$; there are $N_m=4m^3(m-1)$ of them. The links of such a
segment have base points $r_j+te_\nu$ whose coordinates lie in $r_s+[0,
(m-1)d]\subset r_s+[0,md-1)$, so they lie in $\Pi_s$. Write $W_{C}=(W_c)_{c\in
C}$ for a set $C$ of segments and $\mathrm{mid}(C)=\bigcup_{c\in C}
\mathrm{mid}(c)$.

\begin{definition}\label{rC12:C}
\begin{enumerate}[label=\textup{(\roman*)},leftmargin=*]
\item $(\mathrm C_{d,m}(\epsilon))$: for every $s$ and every $\omega\in
G^{E\setminus\Pi_s}$ the law of $W_{C_s}$ under $\nu^\omega_s$ has a density
with respect to $h$ on $G^{C_s}$ with values in $[e^{-\epsilon},e^\epsilon]$.
\item $(\mathrm C^*_d(\eta))$: for every set $C$ of segments and every
configuration of the links outside $\mathrm{mid}(C)$, under the law of
$U_{\mathrm{mid}(C)}$ with density proportional to $\prod_{p\cap
\mathrm{mid}(C)\ne\emptyset}\varphi(U_p)$ the $W_c$, $c\in C$, are
independent and each has a density with respect to $h$ with values in
$[1-\eta,1+\eta]$.
\end{enumerate}
\end{definition}

At $\beta=0$ the law of $W_{C_s}$ under $\nu^\omega_s$ is $h$ for every
$\omega$ (Lemma~\ref{rC12:lines}(iv)), so (i) says that the coarse links
of a patch have a density ratio within $e^{\pm\epsilon}$ with respect to
their $\beta=0$ law, uniformly in the boundary condition. The coarse
plaquettes, and every function of the coarse links, inherit the bound,
since the density of an image law is the conditional expectation of the
density. The law in (ii) is the conditional law under $\mu$ of
$U_{\mathrm{mid}(C)}$ given the links outside $\mathrm{mid}(C)$, and also
under $\nu^\omega_s$ when $\mathrm{mid}(C)\subset\Pi_s$. Condition (ii)
concerns only the local conditional laws of the measure, so it holds for
one measure with these conditional laws exactly when it holds for every
such measure.

\begin{lemma}\label{rC12:mixture}
Assume $(\mathrm C^*_d(\eta))$ with $\eta<1$, let $C$ be a set of segments
and $\mathcal A$ a $\sigma$-algebra contained in $\sigma(U_{E\setminus
\mathrm{mid}(C)})$. Under $\mu$, the conditional law of $W_C$ given
$\mathcal A$ has a density with respect to $h$ with values in
$[(1-\eta)^{\lvert C\rvert},(1+\eta)^{\lvert C\rvert}]$. The same holds
under $\nu^\omega_s$ if $\mathrm{mid}(C)\subset\Pi_s$ and $\mathcal A\subset
\sigma(U_{\Pi_s\setminus\mathrm{mid}(C)})$. In particular
$(\mathrm C^*_d(\eta))$ implies $(\mathrm C_{d,m}(\epsilon))$ with $\epsilon=
N_m\log(1/(1-\eta))$.
\end{lemma}

\begin{proof}
Let $\rho(w)=\prod_{c\in C}\rho_c(w_c)$ be the conditional density of $W_C$
given $U_{E\setminus\mathrm{mid}(C)}$ from $(\mathrm C^*_d(\eta))$; it is
jointly measurable in $w$ and the conditioning links, being given by an
explicit integral. For bounded measurable $F$, by the tower property and
Fubini's theorem, $\E[F(W_C)\mid\mathcal A]=\int F(w)\,\E[\rho(w)\mid\mathcal
A]\,dh(w)$, and $\E[\rho(w)\mid\mathcal A]$ has values in the stated
interval. Under $\nu^\omega_s$ the argument is the same. For the last claim
take $C=C_s$ and $\mathcal A$ trivial under $\nu^\omega_s$, and note
$\log(1+\eta)\le\log(1/(1-\eta))$.
\end{proof}

\begin{lemma}\label{rC12:convolution}
Let $X_1,\dots,X_k$, $k\ge2$, be independent $G$-valued random variables
whose laws have densities $h_i$ with respect to $h$ with $e^{-\delta}\le h_i
\le e^\delta$, $\delta>0$, and let $g,g'\in G$. Then $gX_1\cdots X_kg'$ has
a density $\psi$ with
\[
\sup_G\lvert\psi-1\rvert\le(e^\delta-1)(1-e^{-\delta})^{k-1}.
\]
\end{lemma}

\begin{proof}
For densities $u,v$ the law of the product of independent variables with
these laws has density $u*v(x)=\int u(y)v(y^{-1}x)\,dh(y)=\int u(xz^{-1})
v(z)\,dh(z)$, by invariance of $h$ under translations and inversion. For
integrable $u,v$, $u*1$ and $1*v$ are the constants $\int u$ and $\int v$,
and $u*v$ has integral $(\int u)(\int v)$. Put $u_i=h_i-1$, so $\int u_i=0$.
In the expansion of $(1+u_1)*\cdots*(1+u_k)$ every term that contains both
a factor $1$ and a factor $u_i$ is a convolution of $1$ with a function of
integral zero, hence zero; so $h_1*\cdots*h_k-1=u_1*\cdots*u_k$. Since
$(e^\delta-h_i)(h_i-e^{-\delta})\ge0$ and $\int h_i=1$,
\[
\lVert u_i\rVert_2^2=\int h_i^2-1\le(e^\delta-1)(1-e^{-\delta}).
\]
Write $c=e^{-\delta}$ and $h_i=c+(1-c)g_i$; then $g_i\ge0$ and $\int g_i=1$.
For $v$ with $\int v=0$, $v*u_i=v*h_i=(1-c)v*g_i$, and $v*g_i=\int
v(\cdot\,z^{-1})g_i(z)\,dh(z)$ is an average of right translates of $v$, so
$\lVert v*u_i\rVert_2\le(1-c)\lVert v\rVert_2$ by Minkowski's inequality,
and $v*u_i$ has integral zero. Hence $\lVert u_1*\cdots*u_{k-1}\rVert_2\le
(1-c)^{k-2}\lVert u_1\rVert_2$, and by the Cauchy--Schwarz inequality, for
every $x$,
\[
\begin{split}
\lvert u_1*\cdots*u_k(x)\rvert&\le\lVert u_1*\cdots*u_{k-1}\rVert_2\,\lVert
u_k\rVert_2\\
&\le(1-c)^{k-2}(e^\delta-1)(1-c).
\end{split}
\]
The left side is continuous in $x$, as a convolution of two square
integrable functions. Translation by $g$ and $g'$ preserves $h$, so
$\psi(x)=(h_1*\cdots*h_k)(g^{-1}xg'^{-1})$ obeys the same bound.
\end{proof}

\begin{theorem}\label{rC12:segment}
Let $d\ge4$ and $\delta=6\omega_\varphi>0$. Then $(\mathrm C^*_d(\eta))$
holds with
\[
\eta=(e^\delta-1)(1-e^{-\delta})^{d-3}\le e^\delta\exp\bigl(-(d-3)
e^{-\delta}\bigr).
\]
Consequently $N_m\eta\le1/4$, and $(\mathrm C_{d,m}(\epsilon))$ holds with
$\epsilon\le2N_m\eta$, as soon as $d\ge3+e^\delta(\delta+\log(4N_m))$. For
$\SU(2)$ with $\lvert\beta_W\rvert\le1/20$, $\eta\le0.823\cdot0.452^{d-3}$.
\end{theorem}

\begin{proof}
By Lemma~\ref{rC12:lines}(i),(ii), with $d\ge3$, no plaquette contains
two links of $\mathrm{mid}(C)$: two middle links of one segment lie on one
line, and a plaquette through a middle link contains no other link of a
coarse line. By Lemma~\ref{rC12:lines}(iii), given the links outside
$\mathrm{mid}(C)$ the $U_e$, $e\in\mathrm{mid}(C)$, are independent with
densities in $[e^{-\delta},e^\delta]$. The segments are disjoint, and $W_c=
U_{c,0}\bigl(\prod_{t=1}^{d-2}U_{c,t}\bigr)U_{c,d-1}$ with the end links
$U_{c,0}$, $U_{c,d-1}$ fixed by the conditioning; so the $W_c$ are
independent, and Lemma~\ref{rC12:convolution} with $k=d-2\ge2$ gives the
bound on the density. The second form of $\eta$ uses $e^\delta-1<e^\delta$
and $1-x\le e^{-x}$; with $d-3\ge e^\delta(\delta+\log(4N_m))$ it gives
$\eta\le1/(4N_m)$, and then $\epsilon=N_m\log(1/(1-\eta))\le2N_m\eta$ by
Lemma~\ref{rC12:mixture}, since $\log(1/(1-\eta))\le\eta/(1-\eta)\le2\eta$.
For $\SU(2)$, $\delta=12\lvert\beta_W\rvert\le0.6$, $e^{0.6}-1<0.823$ and
$1-e^{-0.6}<0.452$.
\end{proof}

So the coarse links are close to Haar, in the strong sense of
$(\mathrm C^*_d(\eta))$, uniformly in the boundary condition, for every
compact group, every weight and every coupling, once the cells are large: a
coarse link multiplies $d-2$ conditionally independent group elements, and
products of independent elements with densities bounded below
equidistribute. The proof uses no property of the weight beyond
$\omega_\varphi<\infty$, and the next proposition shows that the
hypothesis does not distinguish a confining from a nonconfining phase.

\begin{proposition}\label{rC12:coulomb}
Let $G=\U(1)$ with the Villain weight $\varphi_\beta$ of
Proposition~\ref{prop:coulomb}, and let $\beta_0$ be as there. For every
$\beta>0$ and every $m\ge4$ there is $d_1$ such that for $d\ge d_1$ the
condition $(\mathrm C^*_d(\eta))$ holds with $N_m\eta\le1/4$, for the
measures of all tori with $d\mid L$ and for the free-boundary measures of
all boxes, with the segments contained in the box. For $\beta>\beta_0$ no
$\mathrm{BP}(n,\varpi)$ holds for the latter measures uniformly in the
boxes of side at least $2n$. Hence $(\mathrm C^*_d(\eta))$, even uniformly
in the volume, does not imply the averaged condition for box measures.
\end{proposition}

\begin{proof}
The Villain weight is a positive continuous class function, so
$\omega_{\varphi_\beta}<\infty$, and Theorem~\ref{rC12:segment} gives
$d_1$ for the tori. Its proof uses only that no plaquette contains two
links of $\mathrm{mid}(C)$ and that a link lies on at most six plaquettes.
Both remain true for the plaquettes of a box with free boundary, and
Lemma~\ref{rC12:lines}(iii) holds there with ``at most six'', so the proof
applies verbatim to box measures. The failure of $\mathrm{BP}$ is
Corollary~\ref{cor:avgcoulomb}.
\end{proof}

Since $(\mathrm C^*_d(\eta))$ concerns only the local conditional laws, it
holds at every $\beta>0$ for every finite-volume measure and every
infinite-volume Gibbs state of the Villain theory, in particular in the
nonconfining phase whose existence Guth proved~\cite{Gut80}; it does not
imply confinement. For tori the paper proves no failure of $\mathrm{BP}$ in
the Coulomb phase, so Proposition~\ref{rC12:coulomb} excludes a proof of
$\mathrm{BP}$ from the coarse-scale hypothesis alone for boxes; a proof for
tori would have to use a property of the torus measures that the box
measures of the same specification lack.

\begin{lemma}\label{rC12:condind}
Let $X$, $Y$ be random variables and $\mathcal A$ a $\sigma$-algebra with
$X$ and $Y$ conditionally independent given $\mathcal A$. Let $X_2$ be a
measurable function of $X$, $Y_2$ one of $Y$, and $\mathcal A'=\mathcal A
\vee\sigma(X_2,Y_2)$. Then $X$ and $Y$ are conditionally independent given
$\mathcal A'$.
\end{lemma}

\begin{proof}
For bounded $\mathcal A$-measurable $a,a'$ and bounded $f(X)$, $g(Y)$,
$\E[afa'g\mid\mathcal A]=aa'\E[f\mid\mathcal A]\E[g\mid\mathcal A]=\E[af
\mid\mathcal A]\E[a'g\mid\mathcal A]$; by the monotone class theorem every
bounded $\sigma(\mathcal A,X)$-measurable function and every bounded
$\sigma(\mathcal A,Y)$-measurable function are conditionally uncorrelated
given $\mathcal A$ in this sense. Fix bounded $f(X)$, $g(Y)$ and put $\phi=
\E[f\mid\mathcal A\vee\sigma(X_2)]$, $\psi=\E[g\mid\mathcal A\vee
\sigma(Y_2)]$; they are $\sigma(\mathcal A,X)$- and $\sigma(\mathcal
A,Y)$-measurable. For bounded $\mathcal A$-measurable $a$ and bounded
$k(X_2)$, $l(Y_2)$,
\[
\begin{split}
\E[fg\,akl]&=\E\bigl[a\,\E[fk\mid\mathcal A]\,\E[gl\mid\mathcal A]\bigr]\\
&=\E\bigl[a\,\E[\phi k\mid\mathcal A]\,\E[\psi l\mid\mathcal A]\bigr]=
\E[\phi\psi\,akl],
\end{split}
\]
and likewise with $g=1$ or $f=1$. The products $akl$ form a class closed
under multiplication that generates $\mathcal A'$, and $\phi$, $\psi$ are
$\mathcal A'$-measurable, so $\E[fg\mid\mathcal A']=\phi\psi$, $\E[f\mid
\mathcal A']=\phi$ and $\E[g\mid\mathcal A']=\psi$.
\end{proof}

\begin{proposition}\label{rC12:nogo}
Let $\varphi$ be nonconstant, $d\ge2$, $K\ge2$, $j\in T_K$, $D=D_j$,
$D'=D_{j+e_0}$, and let $\mathcal O$ be generated by the links outside
$D\cup D'$. Let $y=r_j+(d-1)e_0+e_1+e_2+e_3$, $e=(y,1)$ and $e'=(y+e_0,1)$,
and let $X_2$, $Y_2$ be measurable functions of $U_{D\setminus\{e\}}$ and
of $U_{D'\setminus\{e'\}}$. Then $U_D$ and $U_{D'}$ are not conditionally
independent under $\mu$ given $\mathcal O\vee\sigma(X_2,Y_2)$. In
particular they are not conditionally independent given the links outside
$D\cup D'$ and all the coarse links.
\end{proposition}

\begin{proof}
The base point $y$ lies in $r_j+[0,d)^4$ and $y+e_0=r_{j+e_0}+e_1+e_2+e_3$
lies in $r_{j+e_0}+[0,d)^4$, so $e\in D$ and $e'\in D'$. The coordinate $2$
of $y$ and of $y+e_0$ is $a_2+1\not\equiv a_2$ modulo $d$, so neither link
lies on a coarse line. Hence the coarse links of the segments in $D$ (in
$D'$) are functions of $U_{D\setminus\{e\}}$ (of $U_{D'\setminus\{e'\}}$),
and the other coarse links are functions of the links outside $D\cup D'$;
the last claim follows from the first. Both $e$ and $e'$ lie on the
plaquette $p$ at $y$ in the plane $\{0,1\}$.

If the claim failed, Lemma~\ref{rC12:condind}, applied with $X=U_D$,
$Y=U_{D'}$, $\mathcal A=\mathcal O\vee\sigma(X_2,Y_2)$ and the functions
$U_{D\setminus\{e\}}$, $U_{D'\setminus\{e'\}}$, would make $U_D$ and $U_{D'}$
conditionally independent given all links other than $e$, $e'$, and so
$U_e$ and $U_{e'}$ as well. Their conditional law has a continuous positive
density proportional to $\varphi(AU_{e'}BU_e^{-1})\alpha(U_e)\alpha'(U_{e'})$
with respect to $h\otimes h$. Here $AU_{e'}BU_e^{-1}$, with $A=U_{y,0}$ and
$B=U_{y+e_1,0}^{-1}$, is the plaquette variable of $p$ (replace $\varphi$ by
$\varphi(\cdot^{-1})$ if the orientation is the opposite one), and $\alpha$,
$\alpha'$ collect the other plaquettes through $e$ and through $e'$; no
other plaquette contains both, since two links share at most one
plaquette. Conditional independence for almost every configuration of the
other links gives one configuration at which this continuous density equals
the product of its continuous marginals $h\otimes h$-almost everywhere,
hence everywhere, as $h$ has full support. So $\varphi(AXBY^{-1})=a(X)b(Y)$
for all $X,Y\in G$ with $a,b$ positive, and with $Z=AXB$,
$\varphi(ZY^{-1})=\tilde a(Z)b(Y)$. Setting $Y=1$ and then $Z=1$ gives
$\varphi(ZY^{-1})\varphi(1)=\varphi(Z)\varphi(Y^{-1})$. Thus $\varphi/
\varphi(1)$ is a continuous homomorphism of $G$ into $(\R_{>0},\cdot)$, whose
image is a compact subgroup, hence $\{1\}$, and $\varphi$ is constant, a
contradiction.
\end{proof}

At $\beta=0$ the cells are independent given the coarse links
(Lemma~\ref{rC12:lines}(iv)); the coarse links are coordinates in which the
reference law is a product. For a nonconstant weight every plaquette that
crosses a face between two cells couples them, and by
Proposition~\ref{rC12:nogo} no conditioning on variables of the two cells
that leaves out one pair of links on such a plaquette removes the coupling.
A decoupling of the cells can therefore only be quantitative.

\subsubsection{A quantitative decoupling of the cells}

For $j\in S_s$ let $D^s_j=D_j\cap\Pi_s$; these sets partition $\Pi_s$,
because $r_s+[0,md-1)^4\subset r_s+[0,md)^4=\bigcup_{j\in S_s}(r_j+
[0,d)^4)$. Every segment of $C_s$ starting at $j$ lies in $D^s_j$.

\begin{definition}\label{rC12:T}
$(\mathrm T_{d,m}(\kappa))$: for every $s$, every $\omega\in G^{E\setminus
\Pi_s}$ and every real $g\in L^2(\nu^\omega_s)$, with $\nu=\nu^\omega_s$,
\[
\begin{split}
\nu\bigl(\Var_\nu(g\mid W_{C_s})\bigr)&\le\kappa\sum_{j\in S_s}\nu\bigl(
\Var_\nu(g\mid\mathcal C_j)\bigr),\\
\mathcal C_j&=\sigma\bigl(W_{C_s},U_{\Pi_s\setminus D^s_j}\bigr).
\end{split}
\]
\end{definition}

As for the patch gap, for $L\ge md+2$ this condition depends on $G$,
$\varphi$, $d$, $m$ and $\kappa$ only. At $\beta=0$ the pairs $(U_{D^s_j},
(W_c)_{c\subset D^s_j})$ are independent under $\nu$, so given $W_{C_s}$ the
$U_{D^s_j}$ are independent (as in Lemma~\ref{rC12:lines}(iv)), and the
Efron--Stein inequality for the conditional law gives
$(\mathrm T_{d,m}(1))$. For a nonconstant weight the cells are not
conditionally independent (Proposition~\ref{rC12:nogo}), and
$(\mathrm T_{d,m}(\kappa))$ asks that they be so up to the factor
$\kappa$ in the sense of variances.

\begin{theorem}\label{rC12:reduction}
Let $d\ge4$, $m\ge4$, $d\mid L$, $L\ge2md$. Suppose that
$(\mathrm C^*_d(\eta))$ holds with $N_m\eta\le1/4$, $N_m=4m^3(m-1)$, and
that $(\mathrm T_{d,m}(\kappa))$ holds. Then $\gamma_{d,m}\ge1/(32\kappa+
512)$. If moreover $L$ is even and $m-1\ge148(32\kappa+512)$, then
$\mathrm{BP}(2d-1,1/(1024\kappa+16384))$ holds on $\Lambda_L$, together with
its local and, for $G=\SU(N)$, neutral forms.
\end{theorem}

\begin{proof}
Fix $s$ and $\omega$; write $\nu$, $\Pi$, $\mathcal P_k=\mathcal P^\omega_k$
for $k\in S^\circ_s$, $\mathcal E(g)=\sum_{k\in S^\circ_s}\lVert g-\mathcal
P_kg\rVert^2$ with norms in $L^2(\nu)$, $\mathcal W=\sigma(W_{C_s})$, and
$\mathcal W_{-c}=\sigma(W_{c'}:c'\in C_s,\,c'\ne c)$. Let $g$ be real and
$F=\E_\nu[g\mid\mathcal W]$, so that $\Var g=\nu(\Var(g\mid\mathcal W))+
\Var F$.

\emph{Step 1: cells and blocks.} For $j\in S_s$ let $k(j)\in S^\circ_s$ have
coordinates $k_i=j_i$ if $j_i-s_i\le m-2$ and $k_i=j_i-1$ if $j_i-s_i=m-1$.
Then $D^s_j\subset B_{k(j)}$: in a coordinate with $k_i=j_i$ the offsets
$[0,d)$ lie in $[0,2d-1)$, and with $k_i=j_i-1$ the base points of $D^s_j$
have offsets in $(m-1)d+[0,d-1)$ from $r_s$, inside $(m-2)d+[0,2d-1)$. Each
$k$ is $k(j)$ for at most $16$ cells $j$, and each segment of $C_s$ lies in
the cell of its starting point, so in $B_{k(j)}$; each block receives at
most $64$ segments in this way.

\emph{Step 2: the fine part.} By $(\mathrm T_{d,m}(\kappa))$, and since
$\nu(\Var(g\mid\mathcal A))=\lVert g-\E[g\mid\mathcal A]\rVert^2$ decreases
when $\mathcal A$ grows,
\[
\begin{split}
\nu(\Var(g\mid\mathcal W))&\le\kappa\sum_{j\in S_s}\lVert g-\E[g\mid
U_{\Pi\setminus D^s_j}]\rVert^2\\
&\le\kappa\sum_{j\in S_s}\lVert g-\mathcal
P_{k(j)}g\rVert^2\le16\kappa\,\mathcal E(g),
\end{split}
\]
where the middle step holds because $\mathcal P_{k(j)}g$ is a function of
$U_{\Pi\setminus B_{k(j)}}$, hence of $U_{\Pi\setminus D^s_j}$.

\emph{Step 3: the coarse part.} By Lemma~\ref{rC12:mixture} the law
$\lambda$ of $W_{C_s}$ under $\nu$ has density in $[(1-\eta)^{N_m},
(1+\eta)^{N_m}]$ with respect to $h$. Write $F=\Phi(W_{C_s})$ and let
$Q_c$ integrate out $w_c$ under $h$. By Lemma~\ref{rC12:anova}(ii), and
since $h((\Phi-Q_c\Phi)^2)$ is the least $h((\Phi-\Psi)^2)$ over functions
$\Psi$ of $w_{-c}$,
\[
\begin{split}
\Var_\lambda\Phi&\le(1+\eta)^{N_m}\Var_h\Phi\le(1+\eta)^{N_m}\sum_ch\bigl(
(\Phi-Q_c\Phi)^2\bigr)\\
&\le R^{N_m}\sum_{c\in C_s}\bigl\lVert F-\E[F\mid
\mathcal W_{-c}]\bigr\rVert^2,
\end{split}
\]
with $R=(1+\eta)/(1-\eta)$. As $N_m\ge768$, $\eta\le1/3072$ and $N_m\log R
\le2N_m\eta/(1-\eta)<0.51$, so $R^{N_m}\le2$. Moreover $\E[F\mid\mathcal
W_{-c}]=\E[g\mid\mathcal W_{-c}]$.

\emph{Step 4: leakage.} Let $c\in C_s$ lie in $B_b$, $b=k(j)$ as in Step 1,
let $\mathcal U_b=\sigma(U_{\Pi\setminus B_b})$ and $v=\mathcal P_bg$. The
operator $E_{\mathcal W}-E_{\mathcal W_{-c}}$, with $E_{\mathcal A}$ the
conditional expectation given $\mathcal A$, is an orthogonal projection,
since $\mathcal W_{-c}\subset\mathcal W$; so
\[
\lVert F-\E[g\mid\mathcal W_{-c}]\rVert\le\lVert g-\mathcal P_bg\rVert+
\lVert\psi\rVert,
\]
where $\psi=(E_{\mathcal W}-E_{\mathcal W_{-c}})v$.
Put $\mathcal K=\mathcal U_b\vee\mathcal W_{-c}$. Since $\mathrm{mid}(c)
\subset B_b$ and the segments are disjoint, $\mathcal K\subset\sigma(U_{\Pi
\setminus\mathrm{mid}(c)})$, and by Lemma~\ref{rC12:mixture} the
conditional laws of $W_c$ given $\mathcal K$ and given $\mathcal W_{-c}$
have densities $\rho_{\mathcal K}$, $\rho'$ in $[1-\eta,1+\eta]$. Write
$\psi=\Psi(W_c,W_{-c})$; then $\E[\psi\mid\mathcal W_{-c}]=\int\Psi(w,
W_{-c})\rho'(w)\,dh(w)=0$, and as $W_{-c}$ is $\mathcal K$-measurable,
\[
\begin{split}
\lvert\E[\psi\mid\mathcal K]\rvert&=\Bigl\lvert\int\Psi(w,W_{-c})(\rho_{
\mathcal K}-\rho')(w)\,dh(w)\Bigr\rvert\\
&\le\theta\,\E\bigl[\lvert\psi\rvert
\bigm|\mathcal W_{-c}\bigr],\qquad\theta=\frac{2\eta}{1-\eta},
\end{split}
\]
since $\lvert\rho_{\mathcal K}-\rho'\rvert\le2\eta\le\theta\rho'$. The
function $v$ is $\mathcal K$-measurable and $\nu(\psi)=0$, so
\[
\lVert\psi\rVert^2=\ip\psi v=\ip{\E[\psi\mid\mathcal K]}{v-\nu(v)}\le\theta
\lVert\psi\rVert\,\Var(v)^{1/2},
\]
and $\Var v\le\Var g$. Hence $\lVert F-\E[g\mid\mathcal W_{-c}]\rVert^2\le2
\lVert g-\mathcal P_bg\rVert^2+2\theta^2\Var g$.

\emph{Step 5: conclusion.} Summing over $c$ with Step 1,
$\Var F\le2(128\,\mathcal E(g)+2N_m\theta^2\Var g)$, and with Step 2,
\[
\Var g\le(16\kappa+256)\,\mathcal E(g)+4N_m\theta^2\Var g .
\]
Since $N_m\eta\le1/4$ and $\eta\le1/3072$, $4N_m\theta^2=16N_m\eta^2/(1-
\eta)^2\le4\eta/(1-\eta)^2<1/2$, so $\Var g\le(32\kappa+512)\mathcal E(g)$.
The rest is Theorem~\ref{rC12:finite} with $\gamma=1/(32\kappa+512)$ and
$\varepsilon/16\ge\gamma/32$.
\end{proof}

At $\beta=0$ both hypotheses hold with $\eta=0$, $\kappa=1$, and the theorem
gives $\gamma_{d,m}\ge1/544$, against $\gamma_{d,m}\ge1$ from
Lemma~\ref{rC12:anova}. For $\kappa=1$ the side of the patch is $m\ge80513$
cells, and $(\mathrm C^*_d(\eta))$ with $N_m\eta\le1/4$ then holds once
$d\ge3+e^{6\omega_\varphi}(6\omega_\varphi+\log(4N_m))$ by
Theorem~\ref{rC12:segment}.

\subsubsection{Strong coupling}

\begin{corollary}\label{rC12:strong}
For every $d\ge2$ and $m\ge4$, $\gamma_{d,m}\ge e^{-6(md)^4\omega_\varphi}$.
Hence, for even $L$ with $d\mid L$ and $L\ge594d$, if
\[
\omega_\varphi\le\frac{\log2}{6(297d)^4},
\]
then $\mathrm{BP}(2d-1,1/64)$ holds on $\Lambda_L$, with its local and, for
$G=\SU(N)$, neutral forms. For $\SU(2)$ and $d=2$ (cubes of side three) the
condition is $\lvert\beta_W\rvert\le\log2/(12\cdot594^4)$, which is
$4.6\cdot10^{-13}$ to two digits.
\end{corollary}

\begin{proof}
This is the comparison argument of Holley and Stroock~\cite{HS87}. Let
$\rho$ be the density of $\nu=\nu^\omega_s$ with respect to $h$ on
$G^{\Pi_s}$. A plaquette meeting $\Pi_s$ has its base point in $r_s+
[-1,md-1)^4$, so there are at most $6(md)^4$ of them, and $\rho$ is a
normalised product of their weights; hence $A/a\le e^{6(md)^4\omega_\varphi}$
for $A=\max\rho\ge1\ge a=\min\rho$. Let $Q_j$ integrate out $U_{B_j}$ under
$h$. By Lemma~\ref{rC12:anova}(i), since the $B_j$, $j\in S^\circ_s$,
cover $\Pi_s$, and since $\mathcal P^\omega_jg$ is a function of
$U_{\Pi_s\setminus B_j}$,
\[
\begin{split}
\Var_\nu g&\le A\Var_hg\le A\sum_jh\bigl((g-Q_jg)^2\bigr)\\
&\le A\sum_jh\bigl(
(g-\mathcal P^\omega_jg)^2\bigr)\le\frac Aa\sum_j\nu\bigl((g-\mathcal
P^\omega_jg)^2\bigr).
\end{split}
\]
With $m=297$ and $\gamma=1/2$, $m-1=296\ge148/\gamma$, and
Theorem~\ref{rC12:finite} gives the constant $\gamma/32=1/64$; for
$\SU(2)$, $\omega_\varphi=2\lvert\beta_W\rvert$, and $\log2/(12\cdot594^4)=
4.63979\ldots\cdot10^{-13}$.
\end{proof}

The range is more than eleven orders of magnitude below the range
$\lvert\beta_W\rvert<2/9$ of Corollary~\ref{cor:bpranges}: the comparison
with $\beta=0$ pays for every plaquette of a patch of $(md)^4$ sites, and the
finite-size criterion needs $m$ of order $10^2$. The argument uses neither
Dobrushin coefficients nor single-link estimates.

\begin{remark}[what the hypotheses give]\label{rC12:meaning}
(a) \emph{Proved implications.} $(\mathrm C^*_d(\eta))$ with $N_m\eta\le1/4$
and $(\mathrm T_{d,m}(\kappa))$ give the patch gap
(Theorem~\ref{rC12:reduction}); the patch gap gives
$\mathrm{BP}(2d-1,\varpi)$ on every torus with even $L$, $d\mid L$,
$L\ge2md$ (Theorem~\ref{rC12:finite}), and then decay of correlations by
Theorem~\ref{thm:avgdecay}. Consequently, if along a tuning $(\beta_k,a_k)$
the two hypotheses held with $m$ and $\kappa$ independent of $k$ and
$(2d_k-1)a_k\le\ell$, Theorem~\ref{thm:avgtransfer}, applied with the even
$L\ge2md_k$ divisible by $d_k$, would give every continuum limit of the
torus limits the spectral bound $\sigma(H)\subset\{0\}\cup[m_*,\infty)$
with $m_*=16\varpi/(625e^2\ell)$, $\varpi=1/(1024\kappa+16384)$. The
hypotheses of this implication are not proved along any tuning.

(b) \emph{The coarse-scale hypothesis.} By Theorem~\ref{rC12:segment} and
Proposition~\ref{rC12:coulomb}, $(\mathrm C_{d,m}(\epsilon))$ and
$(\mathrm C^*_d(\eta))$ hold at every coupling for $d$ large compared with
$e^{6\omega_\varphi}$, also in a nonconfining phase, and they do not imply
the averaged condition. Along the weak-coupling tuning the bound of
Theorem~\ref{rC12:segment} does not suffice: for $\SU(2)$ at one loop
$\beta_W\approx4b_0N\log(1/(a\Lambda))$ with $b_0=11N/(48\pi^2)$, which is
$0.3715\log(1/(a\Lambda))$, so $e^{6\omega_\varphi}=e^{12\beta_W}$ grows like
$(a\Lambda)^{-4.46}$, while $d\le\ell/a$. At large $\beta_W$ the
conditional law of a single link is concentrated at scale
$\beta_W^{-1/2}$, and a product of $k$ independent elements with steps of
that size equidistributes after a number of steps of order $\beta_W$. A
single-link contraction of that strength, a factor $1-c/\beta_W$ in place of
the Doeblin factor $1-e^{-6\omega_\varphi}$ in
Lemma~\ref{rC12:convolution}, would give $(\mathrm C^*_d(\eta))$ along
the tuning, since there $d$ may be taken of order $\ell/a$, much larger
than $\beta_W\log\beta_W$. Such a contraction is not proved here, and it would
concern the coarse links only, not the decoupling of the cells.

(c) \emph{The decoupling hypothesis.} $(\mathrm T_{d,m}(\kappa))$ is a
uniform tensorisation of the variance over the cells of side $d$ given the
coarse links, a statement of the same kind as the averaged condition
itself, at the scale of the cells and relative to the coarse links. It holds
with $\kappa=1$ at $\beta=0$. Along the tuning the cells contain $d^4\to
\infty$ sites, and $(\mathrm T_{d,m}(\kappa))$ with $\kappa$ independent of
$k$ is a variance bound at the physical scale $(2d_k-1)a_k\le\ell$. Whether
it fails in the $\U(1)$ Coulomb phase is not decided here: a box version of
Theorem~\ref{rC12:reduction}, together with Theorem~\ref{rC12:segment} and
Corollary~\ref{cor:avgcoulomb}, would show that it fails there for box
measures.

(d) \emph{Relation to $D(n,\alpha)$, (E), the area law and the mass gap.}
The patch gap and $(\mathrm T_{d,m}(\kappa))$ are variance conditions;
$D(n,\alpha)$ is a coupling condition on influences in the forest gauge,
which gives $\mathrm{BP}$ through Proposition~\ref{prop:avgsup}(i) in its
supremum form $D^0_L(n,\alpha)$. No implication between $D(n,\alpha)$ and
$(\mathrm T_{d,m}(\kappa))$ or the patch gap is proved in either direction,
and whether an area law implies $(\mathrm T_{d,m}(\kappa))$ is open.
Nothing in this subsection proves $\mathrm{BP}(n_k,\varpi)$ along the
tuning, $D(n,\alpha)$, condition (E), an area law or the mass gap.
\end{remark}

\begin{remark}[open problems]\label{rC12:open}
(a) $(\mathrm T_{d,m}(\kappa))$, or the patch gap, at any weak coupling
with $d$ large. (b) $(\mathrm C^*_d(\eta))$ along the tuning through a
single-link contraction of order $1/\beta_W$. (c) The box versions of
Theorems~\ref{rC12:finite} and~\ref{rC12:reduction}. (d) A Dobrushin
bound for the patch gap at strong coupling with a range comparable to that
of Corollary~\ref{cor:bpranges}. (e) Better constants: $\kappa=1$ requires
$m\ge80513$, which comes from the factor $p^3(p-2)$ in
Lemma~\ref{rC12:knabe} and the $80$ neighbouring blocks. (f) Whether
$\mathrm{BP}(n,\varpi)$, which is an average over the torus, implies a
patch gap uniform in the boundary condition.
\end{remark}
